\section{Requirements Analysis and Design Methodology}\label{sec:req}

\subsection{Requirements traceability}\label{sec:req-trace}
\Cref{tab:trace} lists the requirements of the brief \cite{tvibrief} with our design answer. ``Implemented''
means configured in the startup-configs and verified as described in \cref{sec:testing}; ``design'' means
specified in this report but not executable on the lab images.

\begin{table}[H]\centering\footnotesize
\caption{Requirements traceability}\label{tab:trace}
\begin{tabularx}{\textwidth}{@{}l >{\raggedright\arraybackslash}p{4.4cm} Y l@{}}
\toprule
\textbf{ID} & \textbf{Requirement} & \textbf{Design answer} & \textbf{Status} \\
\midrule
R1 & Three regions: HQ (2{,}000+), Hanoi (1{,}500+), Da Nang (500+, IoT) & One /16 and one /48 per site; HQ campus with collapsed core; branches with two edge routers & implemented \\
R2 & No single point of failure at L2 and L3 & Paired cores (LACP), dual uplinks, HSRP, two edges and two ISPs per site, two tunnels per branch & implemented \\
R3 & Direct cloud access from branches & Local default route and NAT at each edge & implemented \\
R4 & VLANs, routing and policy provisioned as code & Data model and generators; Netmiko and Ansible & implemented \\
R5 & 4{,}000+ users, 10{,}000+ IoT, guests; 10{,}000+ users by 2030 & VLSM with reserved growth blocks & implemented \\
R6 & 10/40\,Gbit/s backbone, 99.99\,\% uptime & Redundant pairs; 1\,Gbit/s virtual links stand in for 10/40\,Gbit/s & design + lab \\
R7 & Voice, video and data with QoS; VoIP prioritised on the WAN & Marking at the edge, LLQ inside a shaper & implemented \\
R8 & FHRP for gateways, LACP on core links & HSRPv2 on every gateway; Po1 & implemented \\
R9 & VLSM and IPv6; central DHCP with relay & Dual-stack plan; relay to the HQ server & implemented \\
R10 & OSPF inside a site; BGP between sites and to ISPs & OSPFv2/v3 per site; eBGP; iBGP with route reflectors & implemented \\
R11 & Micro-segmentation, NGFW, ACLs for IoT/Guest/Corporate & Per-VLAN ACLs; zone-based firewall at the edge & implemented \\
R12 & Port security, DHCP snooping, DAI & On every access port & implemented \\
R13 & Spine-leaf data centre with EVPN-VXLAN & Fabric design; address block reserved & design \\
R14 & SD-WAN over dual ISP and MPLS & Manual overlay implements the mechanisms; controller design & design \\
R15 & ZTNA / SASE for the hybrid workforce & Target architecture & design \\
R16 & WLC with CAPWAP, Wi-Fi 6/7 & Wireless design & design \\
\bottomrule
\end{tabularx}
\end{table}

\subsection{Assumptions and lab substitutions}\label{sec:req-assumptions}
The lab runs on \todo{EVE-NG / GNS3 version, host CPU and RAM} with Cisco IOSv for routers and IOSvL2 for
switches. Where the lab cannot show a production element, we substitute it and state the consequence
(\cref{tab:substitutions}). All public addresses are documentation prefixes (RFC~5737 \cite{rfc5737},
RFC~3849 \cite{rfc3849}) and TVI uses the private AS~\FactAsTvi{}. Nothing internal is advertised to the
providers, because all Internet traffic is translated.

\begin{table}[H]\centering\footnotesize
\caption{Lab substitutions}\label{tab:substitutions}
\begin{tabularx}{\textwidth}{@{}>{\raggedright\arraybackslash}p{4.3cm} >{\raggedright\arraybackslash}p{3.6cm} Y@{}}
\toprule
\textbf{In the lab} & \textbf{Stands for} & \textbf{Consequence} \\
\midrule
IOSv router \dev{HCMC-DC-SRV01} & DHCP, DNS and NTP servers & Same protocols; server redundancy not shown \\
Zone-based firewall on the edge routers & NGFW pair & Stateful inspection; no application control or IPS \\
1\,Gbit/s virtual links & 10/40\,Gbit/s backbone & Throughput target cannot be measured \\
Shaper at \FactWanRateMbps\,Mbit/s & Contracted ISP rate & Makes WAN congestion reproducible \\
IKEv1 with pre-shared keys, local AAA & IKEv2 with certificates, TACACS+/RADIUS & Simpler lab; listed in the roadmap \\
VPCS hosts & PCs, phones, sensors & Only ping, traceroute and DHCP \\
\bottomrule
\end{tabularx}
\end{table}

\subsection{Design methodology}\label{sec:req-method}
We followed a top-down process in the order of the ten implementation steps of the brief: requirements,
logical topology, addressing, Layer~2, underlay, overlay, services, security and QoS, automation, and
verification. Two working rules shaped the implementation.

\textbf{One source of truth.} Sites, VLANs, subnets, devices, links and tunnels are defined once, in a Python
data model. Generators render from it the startup-configs, the IP plan (Markdown and Excel), the diagrams,
the automation inventory and the tables of this report. Both ends of a link, a BGP session or a tunnel are
therefore produced from the same line of data and cannot disagree, and the report cannot show an address
that is not configured.

\textbf{Verification before the lab.} Errors are cheapest before sixteen virtual devices are cabled. The
configurations went through three layers of checks before deployment: tests of the data model (overlaps,
capacity, unique addresses), tests that read only the rendered configuration files and check that both
ends of every adjacency agree, and an emulation that runs the routing sections with real OSPF and BGP
daemons and injects failures (\cref{sec:test-method}). This process found several design errors that a
syntax check would not have shown; they are discussed in \cref{sec:challenges}.
